\documentclass[11pt,a4paper]{article}
\usepackage[utf8]{inputenc}
\usepackage[T1]{fontenc}
\usepackage[french]{babel}
\usepackage{amsmath,amssymb,amsthm}
\usepackage{geometry}
\geometry{hmargin=2.5cm,vmargin=2.5cm}
\usepackage{fancyhdr}
\usepackage{hyperref}

\pagestyle{fancy}
\fancyhf{}
\lhead{Rapport LaTeX Certifié}
\rhead{\today}
\cfoot{\thepage}

\title{\textbf{Rapport LaTeX Certifié -- Standard Professionnel}}
\author{NeoHarness AI}
\date{\today}

\begin{document}

\maketitle

\begin{abstract}
Ce document présente le rapport LaTeX certifié, rédigé conformément aux exigences de rigueur formelle, typographique et structurelle. Il sert de référence pour la production automatisée de documents scientifiques et techniques de haute qualité.
\end{abstract}

\section{Introduction}
Un rapport LaTeX certifié garantit une mise en page rigoureuse et une conformité typographique irréprochable. Ce standard est particulièrement adapté aux publications scientifiques, aux dossiers techniques et aux rapports institutionnels.

\section{Spécifications Techniques}
Le document respecte les critères suivants :
\begin{itemize}
    \item \textbf{Encodage et polices :} Utilisation native de l'UTF-8 et des fontes vectorielles pour un rendu net sur tous les supports.
    \item \textbf{Environnements mathématiques :} Intégration complète des packages \texttt{amsmath} et \texttt{amssymb}.
    \item \textbf{Structure modulaire :} Séparation claire des sections, sous-sections et annexes.
\end{itemize}

Par exemple, une équation fondamentale s'écrit de manière élégante :
\begin{equation}
    \int_{-\infty}^{+\infty} e^{-x^2} \, dx = \sqrt{\pi}
\end{equation}

\section{Conclusion}
La production de ce rapport atteste de la capacité à générer des livrables documentaires structurés, vérifiés et prêts à l'impression ou à la diffusion numérique.

\end{document}
